\documentclass[11pt]{article}

\usepackage[margin=1in]{geometry}
\usepackage{amsmath,amssymb,amsthm}
\usepackage{graphicx}
\usepackage{hyperref}
\hypersetup{hidelinks}

\newtheorem{theorem}{Theorem}
\newtheorem{lemma}[theorem]{Lemma}
\newtheorem{remark}[theorem]{Remark}
\newcommand{\Pp}{\mathbb P}
\newcommand{\R}{\mathbb R}
\newcommand{\dto}{\mathrel{\mathop{\longrightarrow}\limits^{d}}}

\title{Two obstructions to a uniform weighted stable-limit extension}
\author{Run \texttt{fa\_banach\_001}, lane 6}
\date{August 13, 2026}

\begin{document}
\maketitle

\begin{abstract}
Berkes and Tichy ask whether their uniform weighted subsequence theorem remains
valid, with centering, when the limit random measure merely belongs to the
attraction domain of a fixed stable law.  Two different broad readings are
false.  First, if the fixed stable law may be nonsymmetric, alternating blocks
of signs make the skewness oscillate even when the variables themselves are
strictly stable.  Second, under the broader symmetric regularly-varying-tail
reading, the condition
\[
 \max_{k\leq N}|a_k|=o\!\left(\sum_{k\leq N}|a_k|^\alpha\right)^{1/\alpha}
\]
does not control the true tail scale.  For a symmetric law with
$\Pp(|X|>x)=1/(x\log x)$, we construct positive coefficient blocks satisfying
the displayed condition such that infinitely many endpoint sums, divided by
their only possible nondegenerate scale, converge to the original nonstable
law $X$.  Consequently no normalization and no centering can yield a stable
limit.  Both examples are iid, so passing to a subsequence cannot help.
\end{abstract}

\section{Source question and exact scope}

The source is I. Berkes and R. Tichy, \emph{Lacunary series and stable
distributions}, arXiv:1707.08890 \cite{Source}.  Its Theorem 1.2 produces one
subsequence that works for every real coefficient sequence satisfying the
uniform asymptotic negligibility condition.  Immediately afterward the
authors ask whether the conclusion persists, with a suitable centering
factor, when the limit random measure lies in the domain of normal attraction
of a fixed stable distribution:

\begin{center}
\includegraphics[width=0.96\textwidth]{figures/open_question_crop.png}
\end{center}

There is a terminology boundary that matters.  In strict usage, a normal
domain has a pure $N^{1/\alpha}$ scale (up to a constant), whereas a general
stable attraction domain permits a slowly varying factor.  The source's
uncompiled text after \verb|\end{document}| formulates the intended extension
using balanced regularly varying tails $x^{-\alpha}L(x)$.  We therefore give
two logically independent results:
\begin{enumerate}
\item a counterexample inside a strict normal domain when the fixed stable law
is nonsymmetric; and
\item a symmetric counterexample to the broader regularly-varying-tail
extension, even if the normalization and centering are chosen freely.
\end{enumerate}
We do not claim that the second example refutes the strict symmetric normal-
domain version.

\clearpage
\section{Proof intuition}

The source condition measures coefficients only through their raw
$\ell^\alpha$ energy.  Two extra data become indispensable outside the
symmetric pure-scale setting.

For a nonsymmetric stable law, a positive coefficient contributes one sign of
the stable skewness and a negative coefficient contributes the opposite sign.
Raw absolute energy forgets this information, so very long alternating sign
blocks can force two different endpoint laws.

For a tail $x^{-\alpha}L(x)$ with nonconstant $L$, raw energy also forgets the
scale $L$.  When $L(x)\to0$, many equal small coefficients can carry enormous
raw energy while their random sum remains much smaller.  We exploit that gap
to insert a new coefficient which is negligible in raw energy but dominates
the actual random sum.  At each stage endpoint the weighted sum therefore
looks like one copy of the original, deliberately nonstable, variable.

\section{A convergence-of-types fact}

We use the following standard elementary consequence of tightness.

\begin{lemma}[Convergence of types]\label{lem:types}
Suppose $Y_j\dto Y$, where $Y$ is nondegenerate.  If $r_j>0$, $s_j\in\R$,
and $r_jY_j+s_j\dto Z$ for a nondegenerate $Z$, then every subsequence has a
further subsequence along which
\[
 r_j\longrightarrow r\in(0,\infty),\qquad s_j\longrightarrow s\in\R,
 \qquad Z\stackrel d=rY+s.
\]
In particular, if the transformed sequence has a nondegenerate limit, that
limit is of the same type as $Y$.
\end{lemma}

\begin{proof}
Choose two continuity points of the law of $Y$ whose distribution values are
strictly between zero and one and distinct.  Convergence of the corresponding
quantiles prevents $r_j$ from tending to zero or infinity whenever
$r_jY_j+s_j$ is tight and nondegenerate.  Tightness then also bounds $s_j$.
Extracting a convergent subsequence of $(r_j,s_j)$ and applying Slutsky's
theorem gives $Z\stackrel d=rY+s$.  The same argument applies after starting
from any subsequence.
\end{proof}

\section{The sign obstruction in a strict normal domain}

\begin{theorem}[Nonsymmetric counterexample]\label{thm:sign}
There is an iid determining sequence whose constant limit random measure is a
strictly $1/2$-stable law, and a real coefficient sequence satisfying the
source negligibility condition, for which no positive normalization and
centering can produce one nondegenerate distributional limit.  In particular,
the uniform weighted conclusion is false if the fixed stable law in the
question is allowed to be nonsymmetric.
\end{theorem}

\begin{proof}
Let $X_1,X_2,\ldots$ be iid copies of a positive strictly $1/2$-stable random
variable $X$.  Thus
\[
 X_1+\cdots+X_m\stackrel d=m^2X .
\]
Choose signs $a_k\in\{-1,1\}$ in alternating blocks.  Choose the length of
block $j$ greater than $j$ times the total length of all preceding blocks.
If $N_j$ is its right endpoint and the sign of that block is $\varepsilon_j$,
then the proportion of indices $k\leq N_j$ with sign $\varepsilon_j$ tends to
one.  The signs $\varepsilon_j$ alternate.

For every $N$,
\[
 A_N=\left(\sum_{k=1}^N|a_k|^{1/2}\right)^2=N^2,
 \qquad \frac{\max_{k\leq N}|a_k|}{A_N}=N^{-2}\longrightarrow0.
\]
Let $p_N$ and $q_N$ be the proportions of positive and negative signs among
$a_1,\ldots,a_N$.  Stability and independence give the exact identity
\[
 N^{-2}\sum_{k=1}^Na_kX_k
 \stackrel d=p_N^2X^+-q_N^2X^-,
\]
where $X^+,X^-$ are independent copies of $X$.  Hence at positive block
endpoints the normalized sums converge to $X$, while at negative block
endpoints they converge to $-X$.

Suppose arbitrary centerings and positive normalizers produced one
nondegenerate limit along all $N$.  Lemma~\ref{lem:types}, applied first to the
positive and then to the negative endpoints, would make that limit both a
positive affine transform of $X$ and a positive affine transform of $-X$.
The former is bounded below and unbounded above; the latter is bounded above
and unbounded below.  They cannot have the same nondegenerate law.  Finally,
an iid sequence is determining with constant limit random measure
$\mathcal L(X)$, and every subsequence remains iid, so subsequence selection
cannot remove the obstruction.
\end{proof}

\section{The slowly varying scale obstruction}

We now retain symmetry and use only positive coefficients.

\begin{theorem}[Balanced regular-variation counterexample]\label{thm:slow}
There exist a symmetric law $F$ in the general domain of attraction of a
symmetric $1$-stable law, an iid determining sequence $(X_n)$ with law $F$,
and positive coefficients $(a_k)$ such that
\[
 \frac{\max_{k\leq N}a_k}{\sum_{k=1}^Na_k}\longrightarrow0,
\]
but for no positive normalizers $d_N$ and centerings $c_N$ does
\[
 d_N^{-1}\left(\sum_{k=1}^Na_kX_k-c_N\right)
\]
converge to a nondegenerate stable law.  The same assertion holds after
replacing $(X_n)$ by any subsequence.
\end{theorem}

\begin{proof}
Let $R\geq e$ have survival function
\[
 \Pp(R>x)=\frac1{x\log x}\qquad(x\geq e),
\]
with the remaining probability placed at $e$, and let $X=\epsilon R$, where
$\epsilon$ is an independent symmetric sign.  Then
\[
 \Pp(|X|>x)=\frac1{x\log x}\qquad(x\geq e).
\]
The tails are balanced and regularly varying with exponent $-1$.  The
classical stable domain-of-attraction theorem therefore puts $X$ in the
domain of attraction of a symmetric $1$-stable law.  More explicitly, if
$b_m>e$ is defined by
\[
 b_m\log b_m=m,
\]
then $b_m\sim m/\log m$ and
\[
 b_m^{-1}(X_1+\cdots+X_m)
\]
converges in distribution to a symmetric $1$-stable law; in particular this
sequence is tight.  The law of $X$ itself is not stable: its tail is
$x^{-1}/\log x$, whereas a non-Gaussian stable law has a nonzero constant
multiple of a power tail (and a Gaussian stable law has a light tail).

We construct the coefficients recursively.  Put $w_0=1$.  At stage $j$,
append $m_j$ copies of $w_{j-1}$ and then append one coefficient
\[
 w_j=\frac{w_{j-1}m_j}{\sqrt{\log m_j}}.                 \tag{1}
\]
Let $T_{j-1}$ denote the weighted sum through the preceding stage.  When
$m_j$ is selected, $T_{j-1}$ is already fixed in law, while the right side of
(1) tends to infinity with $m_j$.  We may therefore choose integers
$m_j\to\infty$ so fast that
\[
 \Pp\{|T_{j-1}|>w_j/j\}<1/j,
 \qquad (\log m_j)^{-1/2}<1/j.                          \tag{2}
\]

Let $M_N=\max_{k\leq N}a_k$ and $A_N=\sum_{k\leq N}a_k$.  During the block of
$m_j$ copies, $M_N/A_N$ can only decrease from its value at the preceding
endpoint.  Upon appending $w_j$, it is at most
\[
 \frac{w_j}{m_jw_{j-1}}=\frac1{\sqrt{\log m_j}}.
\]
Induction over the stages and (2) show $M_N/A_N\to0$ over every index $N$.

At the endpoint of stage $j$, write
\[
 T_j=T_{j-1}+w_{j-1}\sum_{r=1}^{m_j}X_{j,r}+w_jX_j^*,
\]
using fresh iid variables.  The first term divided by $w_j$ tends to zero in
probability by (2).  For the middle term, tightness at scale $b_{m_j}$ and
(1) give
\[
 \frac{w_{j-1}b_{m_j}}{w_j}
 \sim\frac1{\sqrt{\log m_j}}\longrightarrow0.
\]
Consequently Slutsky's theorem yields
\[
 \frac{T_j}{w_j}\dto X.                                \tag{3}
\]

If some centerings $c_N$ and positive normalizers $d_N$ gave a nondegenerate
stable limit along the full sequence, they would do so along the stage
endpoints.  Apply Lemma~\ref{lem:types} to (3), with
$r_j=w_j/d_{N_j}$ and $s_j=-c_{N_j}/d_{N_j}$.  The alleged stable limit would
have to be an affine transform of $X$, contradicting that $X$ is not stable.
As before, every subsequence of an iid sequence has the same joint laws, so
the single deterministic coefficient construction defeats every selected
subsequence.
\end{proof}

\section{Consequences, limitations, and novelty audit}

Theorem~\ref{thm:sign} shows that ``symmetry plays no essential role'' needs a
qualification when arbitrary real coefficients are allowed: one must at least
control the signed $\alpha$-energy, or allow the target skewness to depend on
its limiting balance.  Theorem~\ref{thm:slow} shows that balanced regular
variation alone is also insufficient under the source coefficient condition.
A viable general-domain theorem needs a negligibility condition in the actual
tail scale, involving $L$ as well as $|a_k|^\alpha$.

The slow-variation example is not in a strict symmetric normal domain because
$L(x)=1/\log x$ does not tend to a positive constant.  Accordingly we do not
claim to settle a version that simultaneously insists on a symmetric fixed
stable law and strict normal-domain terminology.  In that regime the leading
small-characteristic-function exponent is aligned with the hypothesis used
in the source proof, after the appropriate centering.

A bounded search on August 13, 2026 checked all four lightweight run indexes
for the arXiv id, exact title, and the weighted stable-domain keywords.
External exact-title and weighted-sum searches found the source and classical
triangular-array results whose hypotheses explicitly retain the slowly
varying characteristic factor, but no later paper resolving this source
question.  The construction here is explicit and self-contained, but priority
remains subject to specialist review.

The packet is recommended for expert review as a likely valid pair of full
scoped counterexamples.  The principal audit points are the standard
positive-stable scaling in Theorem~\ref{thm:sign}, the recursive quantifiers in
Theorem~\ref{thm:slow}, and the distinction between strict normal attraction
and general regularly varying attraction.

\begin{thebibliography}{9}
\bibitem{Source}
I. Berkes and R. Tichy,
\emph{Lacunary series and stable distributions},
Periodica Mathematica Hungarica 78 (2019), 47--61;
arXiv:1707.08890.

\bibitem{Feller}
W. Feller,
\emph{An Introduction to Probability Theory and Its Applications, Vol. II},
2nd ed., Wiley, 1971, Chapter XVII.
\end{thebibliography}

\end{document}
